\section{Object 2}
\markright{Object 2}

\subsubsection{Part 1}

The first component was designed with a specialized profile that generates motion through geometric interaction with mating elements. Complex curved surfaces were created using sketch and extrude features with precise dimensional control. The part geometry was carefully refined to ensure smooth operation and efficient energy transmission throughout the operational cycle.

\begin{figure}[H]
	\centering
	\begin{subfigure}[b]{0.4\textwidth}
		\centering
		\includegraphics[width=\textwidth, keepaspectratio]{lab-03/assets/2/1.jpg}
	\end{subfigure}%
	\hspace{0.3em}
	\begin{subfigure}[b]{0.4\textwidth}
		\centering
		\includegraphics[width=\textwidth, keepaspectratio]{lab-03/assets/2/2.jpg}
	\end{subfigure}%
	\caption{Snapshots while Designing Part 1 of Object 2 in Solidworks}
	\label{lab03_object_2_part_1}
\end{figure}

\subsubsection{Part 2}

This component was modeled to respond to the geometric inputs provided by the primary part, converting its motion into linear or rotational output. A cylindrical body with appropriate contact surfaces was created using revolve and bore features. The design included specific dimensions and surface characteristics to ensure smooth interaction and minimize wear during operation.

\begin{figure}[H]
	\centering
	\begin{subfigure}[b]{0.4\textwidth}
		\centering
		\includegraphics[width=\textwidth, keepaspectratio]{lab-03/assets/2/3.jpg}
	\end{subfigure}%
	\hspace{0.3em}
	\begin{subfigure}[b]{0.4\textwidth}
		\centering
		\includegraphics[width=\textwidth, keepaspectratio]{lab-03/assets/2/4.jpg}
	\end{subfigure}%
	\caption{Snapshots while Designing Part 2 of Object 2 in Solidworks}
	\label{lab03_object_2_part_2}
\end{figure}

\subsubsection{Part 3}

A structural support element was designed to constrain the motion of the primary components and maintain proper spatial relationships. Extrusion and pocket features were employed to create mounting surfaces and guide slots for other parts. Appropriate clearances were incorporated to allow smooth motion while maintaining rigidity of the overall structure.

\begin{figure}[H]
	\centering
	\begin{subfigure}[b]{0.4\textwidth}
		\centering
		\includegraphics[width=\textwidth, keepaspectratio]{lab-03/assets/2/5.jpg}
	\end{subfigure}%
	\hspace{0.3em}
	\begin{subfigure}[b]{0.4\textwidth}
		\centering
		\includegraphics[width=\textwidth, keepaspectratio]{lab-03/assets/2/6.jpg}
	\end{subfigure}%
	\caption{Snapshots while Designing Part 3 of Object 2 in Solidworks}
	\label{lab03_object_2_part_3}
\end{figure}


\subsubsection{Assembly}

The three components were assembled with coincident and parallel constraints to establish the kinematic linkages required for proper mechanism function. The assembly was analyzed to verify that the motion sequences are correct and no component interference occurs during operation. Motion simulation confirmed that the mechanism achieves the intended output motion as designed.

\begin{figure}[H]
	\centering
	\begin{subfigure}[b]{0.32\textwidth}
		\centering
		\includegraphics[width=\textwidth, keepaspectratio]{lab-03/assets/2/7.jpg}
	\end{subfigure}%
	\hspace{0.3em}
	\begin{subfigure}[b]{0.32\textwidth}
		\centering
		\includegraphics[width=\textwidth, keepaspectratio]{lab-03/assets/2/8.jpg}
	\end{subfigure}%
	\hspace{0.3em}
	\begin{subfigure}[b]{0.32\textwidth}
		\centering
		\includegraphics[width=\textwidth, keepaspectratio]{lab-03/assets/2/9.jpg}
	\end{subfigure}%
	\caption{Snapshots while Assembling Object 2 in Solidworks}
	\label{lab03_object_2_assembly}
\end{figure}

